\chapter{General Organization of the Central Unit}

\begin{khulasa}[: what will we learn in this chapter?]
\begin{itemize}
    \item What a computer is, and how it changed through history.
    \item The main parts: the processor, the memory, the input/output devices.
    \item The Von Neumann and Harvard models, and the difference between them.
    \item The types of memory, and how to calculate a memory size.
    \item The three buses: the data bus, the address bus, the control bus.
\end{itemize}
\end{khulasa}

%================================================
\section{What is a computer?}
%================================================

\begin{tarif}
    A \term{computer}{ordinateur} is an electronic machine. It receives \textbf{data},
    processes it with a \textbf{program} (a list of instructions), and gives \textbf{results}.
    It can also store data and results.
\end{tarif}

Everything a computer does fits in four verbs: it \textbf{receives}, it \textbf{processes},
it \textbf{stores}, it \textbf{outputs}.

\begin{center}
\begin{tikzpicture}[node distance=1.4cm]
    \node[unit=brandGreen, minimum width=2.8cm, minimum height=1.2cm] (in) {Input\\{\small keyboard, mouse}};
    \node[unit, minimum width=2.8cm, minimum height=1.2cm, right=of in] (cpu) {Processing\\{\small the processor}};
    \node[unit=brandOrange, minimum width=2.8cm, minimum height=1.2cm, right=of cpu] (out) {Output\\{\small screen, printer}};
    \node[unit=brandPurple, minimum width=2.8cm, minimum height=1.1cm, below=1cm of cpu] (mem) {Storage\\{\small memory, disk}};
    \draw[sahmT] (in) -- (cpu);
    \draw[sahmT] (cpu) -- (out);
    \draw[sahmT, <->] (cpu) -- (mem);
\end{tikzpicture}
\end{center}

Let us take a simple example: you use the calculator app on your computer to compute $25 \times 4$.
\begin{enumerate}
    \item \textbf{Input}: you type 25, then $\times$, then 4 on the keyboard. Each key press becomes a
          number code that goes into the computer.
    \item \textbf{Temporary storage}: the computer keeps the two numbers in memory, because it needs them soon.
    \item \textbf{Processing}: the processor reads the two numbers, runs the multiply instruction, and gets 100.
    \item \textbf{Output}: the computer sends the result to the screen, and you see it.
\end{enumerate}
Every program you use, even a complex one (a game, a browser, a photo app), is in fact millions of
these small steps, done very fast.

\begin{tashbih}
    A computer is like a \textbf{restaurant kitchen}. The ingredients (vegetables, meat) are the \textbf{data}.
    The recipe is the \textbf{program}. The cook is the \textbf{processor}. The fridge is the \textbf{memory}.
    The finished dish is the \textbf{result}. The cook does not decide alone: he follows the recipe step by step.
\end{tashbih}

\subsection{Hardware and software}
\begin{itemize}
    \item \term{Hardware}{matériel}: everything you can touch: the processor, the memory, the screen, the wires.
    \item \term{Software}{logiciel}: the programs that tell the hardware what to do: the operating
          system, the browser, games.
\end{itemize}
In this course we study \textbf{the hardware from the inside}, and we see how it runs programs.

\subsection{The layers of a computer}
We can see a computer as layers, one on top of the other. Each layer uses the services of the layer
below it, without knowing all its details:

\begin{center}
\begin{tikzpicture}[l/.style={draw=#1, fill=#1!15, very thick, minimum width=9cm, minimum height=0.75cm, font=\small}]
    \node[l=brandGreen] at (0,1.8) {Applications: browser, games, word processor};
    \node[l=brandPurple] at (0,0.9) {Operating system: Windows, Linux, Android};
    \node[l=brandOrange] at (0,0) {Machine language and assembly: what the processor understands};
    \node[l=brandTeal] at (0,-0.9) {Hardware: processor, memory, buses, electronic circuits};
    \draw[sahmT] (5,-1.1) -- (5,2) node[midway, right, font=\small] {from machine to human};
\end{tikzpicture}
\end{center}
In this course we work on the \textbf{two lowest layers}: the hardware, and the assembly language
that talks to it directly.

\subsection{How does a program run?}
When you double-click on a program icon, these steps happen:
\begin{enumerate}
    \item The program is stored on the \textbf{disk} (SSD or HDD) as a file of zeros and ones.
    \item The operating system \textbf{copies it into the main memory} (RAM), because the disk is very slow
          and the processor does not read from it directly.
    \item The processor reads the instructions from memory \textbf{one after the other} and runs them.
    \item When you close the program, its space in memory is freed for other programs.
\end{enumerate}

%================================================
\section{A short history of computers}
%================================================

The computer did not appear in one day. These are the main steps:

\begin{center}
\begin{tikzpicture}[
    yr/.style={circle, fill=brandTeal, minimum size=0.35cm, inner sep=0pt},
    txt/.style={text width=2.7cm, align=center, font=\small}]
    \draw[brandTeal, line width=3pt] (0,0) -- (14,0);
    \foreach \x/\y/\t/\p in {
        0/1642/{Pascal's machine\\ adds, subtracts}/above,
        2/1837/{Babbage's machine\\ idea of a program}/below,
        4/1945/{ENIAC\\ vacuum tubes}/above,
        6/1947/{invention of the\\ \textbf{transistor}}/below,
        8/1958/{\textbf{integrated circuit}\\ (IC)}/above,
        10/1971/{first microprocessor\\ Intel 4004}/below,
        12/1978/{the\\ Intel 8086}/above,
        14/2005/{\textbf{multi-core}\\ processors}/below} {
        \node[yr] (n) at (\x, 0) {};
        \node[txt, \p=0.25cm of n] {\textbf{\y}\\ \t};
    }
\end{tikzpicture}
\end{center}

\begin{table}[h]
\centering
\begin{tabular}{|c|c|c|c|}
    \hline
    \ras{Generation} & \ras{Period} & \ras{Main technology} & \ras{Example} \\ \hline
    1st & 1945--1955 & vacuum tubes & ENIAC (30 tons!) \\ \hline
    2nd & 1955--1965 & transistor & IBM 7090 \\ \hline
    3rd & 1965--1980 & integrated circuits & IBM 360 \\ \hline
    4th & 1980--today & microprocessor (VLSI) & personal computer, phone \\ \hline
\end{tabular}
\end{table}

Each new generation brought a technology that made the parts \textbf{smaller}, \textbf{faster}
and \textbf{less hungry for energy}:
\begin{itemize}
    \item \textbf{Vacuum tubes} were big, got very hot, and burned out all the time. Some had to be changed almost every day.
    \item A \textbf{transistor} does the same job (an electric switch that opens and closes), but it is hundreds of times smaller and does not burn out.
    \item An \textbf{integrated circuit} puts thousands of transistors on one small piece of silicon, instead of connecting them with wires.
    \item A \textbf{microprocessor} puts the \textbf{whole} processor on one chip. This gave us the personal computer, then the smartphone.
\end{itemize}

\begin{tarif}
    \term{Moore's law}{loi de Moore}: in 1965, Gordon Moore (a founder of Intel) noticed that the number
    of transistors on a chip \textbf{doubles about every two years}. This stayed true for more than fifty
    years: from 2\,300 transistors in the 4004 to tens of billions today.
\end{tarif}

\begin{hal}
    ENIAC filled a whole room and used 150 kilowatts. Your phone is millions of times more powerful and runs
    on a small battery. The reason: transistors became tiny, so one processor now holds billions of them.
\end{hal}

%================================================
\section{Quick reminder: how does a computer count?}
%================================================

A computer only knows two electric states: \textbf{current flows} or \textbf{current does not flow}.
We write them 1 and 0.

\begin{tarif}
    \begin{itemize}
        \item A \term{bit}{bit}: the smallest unit of information. Its value is 0 or 1.
        \item A \term{byte}{octet}: a group of \textbf{8 bits}.
        \item A \term{word}{mot}: the number of bits the processor handles at once (16 bits in the 8086).
    \end{itemize}
\end{tarif}

We write numbers in three systems: decimal (base 10), binary (base 2), and hexadecimal (base 16, its
digits are 0 to 9, then A to F).

\begin{center}
\small
\begin{tabular}{|c|c|c||c|c|c|}
    \hline
    \ras{Decimal} & \ras{Binary} & \ras{Hex} & \ras{Decimal} & \ras{Binary} & \ras{Hex} \\ \hline
    0 & 0000 & 0 & 8  & 1000 & 8 \\ \hline
    1 & 0001 & 1 & 9  & 1001 & 9 \\ \hline
    2 & 0010 & 2 & 10 & 1010 & A \\ \hline
    3 & 0011 & 3 & 11 & 1011 & B \\ \hline
    4 & 0100 & 4 & 12 & 1100 & C \\ \hline
    5 & 0101 & 5 & 13 & 1101 & D \\ \hline
    6 & 0110 & 6 & 14 & 1110 & E \\ \hline
    7 & 0111 & 7 & 15 & 1111 & F \\ \hline
\end{tabular}
\end{center}

\begin{khatwa}[: convert \hex{2F} to binary and to decimal]
    \begin{enumerate}
        \item Each hex digit is \textbf{4 bits}: $2 = 0010$ and $F = 1111$.
        \item Put them side by side: $\hex{2F} = \bin{0010\,1111}$.
        \item To decimal: $2 \times 16 + 15 = 32 + 15 = 47$.
    \end{enumerate}
\end{khatwa}

\begin{khatwa}[: convert the binary number 1011 to decimal]
    Each bit has a ``weight'', a power of 2. We start on the right with $2^0 = 1$:
    \begin{center}
    \begin{tabular}{|c|c|c|c|c|}
        \hline \ras{Bit} & 1 & 0 & 1 & 1 \\ \hline
        \ras{Weight} & $2^3 = 8$ & $2^2 = 4$ & $2^1 = 2$ & $2^0 = 1$ \\ \hline
        \ras{Bit $\times$ weight} & 8 & 0 & 2 & 1 \\ \hline
    \end{tabular}
    \end{center}
    We add: $8 + 0 + 2 + 1 = 11$. So $\bin{1011} = 11$.
\end{khatwa}

\begin{khatwa}[: convert the decimal number 13 to binary (divide by 2 again and again)]
    \begin{enumerate}
        \item $13 \div 2 = 6$, remainder \textbf{1}.
        \item $6 \div 2 = 3$, remainder \textbf{0}.
        \item $3 \div 2 = 1$, remainder \textbf{1}.
        \item $1 \div 2 = 0$, remainder \textbf{1}. We stop, because the quotient is 0.
        \item Read the remainders \textbf{from the last to the first}: \bin{1101}. Check: $8 + 4 + 0 + 1 = 13$.
    \end{enumerate}
\end{khatwa}

\begin{tashbih}
    Why use hexadecimal? Because binary numbers are long and tiring to read: \bin{1010\,1111\,0011\,0110}.
    Hexadecimal packs every 4 bits into one digit, so the number becomes \hex{AF36}. It is the
    ``short form'' programmers use to write binary quickly, like writing ``Dr.'' instead of ``Doctor''.
\end{tashbih}

\begin{intabih}
    In assembly language we write \code{h} after a hex number (\code{2Fh}), and \code{b} after a binary
    number (\code{00101111b}). A number without a letter is decimal.
\end{intabih}

\subsection{Memory size units}
\begin{center}
\begin{tabular}{|c|c|c|}
    \hline
    \ras{Unit} & \ras{Symbol} & \ras{Value} \\ \hline
    kilobyte & KB & $2^{10} = 1024$ bytes \\ \hline
    megabyte & MB & $2^{20} = 1\,048\,576$ bytes \\ \hline
    gigabyte & GB & $2^{30}$ bytes (about one billion) \\ \hline
    terabyte & TB & $2^{40}$ bytes \\ \hline
\end{tabular}
\end{center}

%================================================
\section{The main parts of a computer}
%================================================

Inside every computer, from the biggest server to the smallest smartwatch, we find three main parts:

\begin{center}
\begin{tikzpicture}
    \node[unit, minimum width=4cm, minimum height=2.6cm] (cpu) at (0,0) {};
    \node[font=\bfseries] at (0, 1) {Processor (CPU)};
    \node[unit=brandOrange, minimum width=1.6cm, font=\small] at (-0.95, 0) {UC\\control};
    \node[unit=brandGreen, minimum width=1.6cm, font=\small] at (0.95, 0) {UAL\\compute};
    \node[draw=brandDark, fill=white, font=\small, minimum width=3.4cm] at (0,-0.9) {registers};
    \node[unit=brandPurple, minimum width=3.4cm, minimum height=2.6cm] (mem) at (-6,0) {Main memory\\{\small RAM and ROM}};
    \node[unit=brandGreen, minimum width=3.4cm, minimum height=2.6cm] (io) at (6,0) {Input/output\\devices};
    \draw[sahmT, <->] (cpu) -- (mem);
    \draw[sahmT, <->] (cpu) -- (io);
\end{tikzpicture}
\end{center}

\begin{enumerate}
    \item The \term{processor}{processeur} (CPU): the ``brain'' of the computer. It reads the instructions and runs them.
    \item The \term{main memory}{mémoire centrale}: it keeps the program that is running now, and its data.
    \item The \term{input/output devices}{périphériques d'E/S}: they let the computer talk with the outside world.
\end{enumerate}

\begin{khatwa}[: the trip of a letter from the keyboard to the screen]
    You press the key \textbf{A} in a text editor:
    \begin{enumerate}
        \item The keyboard (\textbf{input}) sends the code of the key to the computer.
        \item The processor is interrupted (it stops its work for a moment), reads the code, and stores it in
              \textbf{memory}, inside the text you are writing.
        \item The text editor asks the processor to ``draw'' the letter: the processor computes the light dots that make the letter.
        \item These dots go to the graphics card, then to the screen (\textbf{output}).
    \end{enumerate}
    All this happens in a tiny part of a second, so the letter seems to appear at once.
\end{khatwa}

\begin{tarif}
    The \term{central unit}{unité centrale}: the processor plus the main memory. This is the part that really runs the programs.
\end{tarif}

%================================================
\section{The two basic models: Von Neumann and Harvard}
%================================================

\subsection{The Von Neumann model}

In 1945, the scientist John von Neumann proposed a simple and brilliant idea: we keep the
\textbf{program} and the \textbf{data} together in \textbf{the same memory}. Before that, changing the
program meant re-connecting wires by hand!

In Von Neumann's model, a computer has \textbf{five units}: input, output, memory, the arithmetic and
logic unit, and the control unit. The key idea is the \term{stored program}{programme enregistré}: instructions
are numbers, just like data, so we can keep them in the same memory, and read and change them the same way.

\begin{center}
\begin{tikzpicture}
    \node[unit, minimum width=4.2cm, minimum height=2.4cm] (cpu) at (0,0) {};
    \node[font=\bfseries] at (0,0.85) {Processor};
    \node[unit=brandOrange, font=\small, minimum width=1.7cm] at (-1,-0.15) {control unit};
    \node[unit=brandGreen, font=\small, minimum width=1.7cm] at (1,-0.15) {UAL};
    \node[unit=brandPurple, minimum width=4.2cm, minimum height=2.4cm, align=center] (mem) at (-7,0) {};
    \node[font=\bfseries] at (-7,0.85) {one memory};
    \node[draw=brandDark, fill=brandTeal!20, font=\small, minimum width=3.4cm] at (-7,0.1) {program (instructions)};
    \node[draw=brandDark, fill=brandOrange!20, font=\small, minimum width=3.4cm] at (-7,-0.6) {data};
    \draw[brandDark, line width=5pt, {Stealth[length=4mm]}-{Stealth[length=4mm]}] (-4.8,0) -- (-2.2,0);
    \node[font=\small, above] at (-3.5,0.2) {one bus};
    \node[unit=brandGreen, minimum width=2.6cm, minimum height=1.2cm] (io) at (5,0) {input and\\output};
    \draw[sahmT, <->] (cpu) -- (io);
\end{tikzpicture}
\end{center}

\begin{itemize}
    \item \textbf{Advantage}: simple and flexible. To change the program, we just load another program into memory.
    \item \textbf{Drawback}: one bus for instructions and data, so we cannot read an instruction and a data
          item at the same moment. This traffic jam is called the \term{Von Neumann bottleneck}{goulot de Von Neumann}.
\end{itemize}

\begin{tashbih}
    Imagine a library with \textbf{one narrow door}. The students who bring back books (the data) and the
    students who come for instructions (the program) all use the same door, one by one. Even if the staff
    inside is very fast, the door decides the speed. This is the bottleneck.
\end{tashbih}

\subsection{The Harvard model}
In this model we use \textbf{two separate memories}: one for instructions and one for data. Each one has its own bus.

\begin{center}
\begin{tikzpicture}
    \node[unit, minimum width=3.4cm, minimum height=2.2cm] (cpu) at (0,0) {Processor};
    \node[unit=brandTeal, minimum width=3.6cm, minimum height=1.6cm] (im) at (-6.5,0) {instruction memory\\{\small the program}};
    \node[unit=brandOrange, minimum width=3.6cm, minimum height=1.6cm] (dm) at (6.5,0) {data memory\\{\small the variables}};
    \draw[tealDeep, line width=4pt, -{Stealth[length=4mm]}] (im) -- (cpu) node[midway, above, font=\small, text=brandDark] {instruction bus};
    \draw[orangeDeep, line width=4pt, {Stealth[length=4mm]}-{Stealth[length=4mm]}] (cpu) -- (dm) node[midway, above, font=\small, text=brandDark] {data bus};
\end{tikzpicture}
\end{center}

\begin{itemize}
    \item \textbf{Advantage}: faster, because the processor reads the next instruction and a data item at the same time.
    \item \textbf{Drawback}: more complex and more expensive, and the size of each memory is fixed.
\end{itemize}

\begin{tashbih}
    In the Harvard library there are \textbf{two doors}: one for instructions and one for data. The two
    queues move at the same time, so people get in faster. But two doors cost more, and if one side is full
    we cannot use the free space of the other side.
\end{tashbih}

\subsection{Comparison and uses}
\begin{center}
\begin{tabularx}{\linewidth}{|>{\centering\arraybackslash}p{3cm}|C|C|}
    \hline
    \ras{} & \ras{Von Neumann} & \ras{Harvard} \\ \hline
    Memory   & one memory for program and data & two separate memories \\ \hline
    Buses    & one shared bus & two independent buses \\ \hline
    Speed    & slower (bottleneck) & faster (two reads at once) \\ \hline
    Complexity & simple and cheap & more complex, more expensive \\ \hline
    Used in  & personal computers, servers & microcontrollers (Arduino), signal processors (DSP) \\ \hline
\end{tabularx}
\end{center}

\begin{hal}
    The processors in your computer mix the two models. The main memory (RAM) is shared (Von Neumann), but
    inside the processor there are two small, fast, separate memories: one for instructions and one for data
    (Harvard). We call this \textbf{modified Harvard}. We come back to it in Chapter 5.
\end{hal}

%================================================
\section{The processor: role and parts}
%================================================

\begin{tarif}
    The \term{processor}{processeur, UCT} (CPU): the circuit that runs the instructions of a program, one after the other.
\end{tarif}

The processor has three main parts (we study them in detail in Chapter 2):

\begin{center}
\begin{tabularx}{\linewidth}{|>{\centering\arraybackslash}p{4.2cm}|R|R|}
    \hline
    \ras{Part} & \ras{Its job} & \ras{Analogy} \\ \hline
    control unit (UC) & reads the instruction, understands it, gives orders to the other parts & the orchestra conductor \\ \hline
    arithmetic and logic unit (UAL) & does arithmetic ($+$, $-$) and logic (AND, OR) & the calculator \\ \hline
    registers & very small and very fast memory inside the processor & the scrap paper on the desk \\ \hline
\end{tabularx}
\end{center}

\subsection{The clock and the speed of the processor}
The processor works to the rhythm of a \term{clock}{horloge} that sends regular pulses. At each pulse the
processor moves one small step forward in its work.

\begin{tarif}
    The \textbf{frequency} $f$: the number of pulses per second, measured in hertz (Hz).
    The \textbf{period} $T$: the length of one pulse, and $T = \dfrac{1}{f}$.
\end{tarif}

\begin{khatwa}[: a processor at 3 GHz]
    \begin{enumerate}
        \item $f = 3\ \text{GHz} = 3 \times 10^{9}$ pulses per second.
        \item $T = \dfrac{1}{3 \times 10^{9}} \approx 0.33 \times 10^{-9}$ second, so about 0.33 nanosecond.
        \item In the time you blink (about 0.1 second), this processor does 300 million cycles!
    \end{enumerate}
\end{khatwa}

\begin{intabih}
    The frequency alone does not decide the speed of a processor. What matters is the \textbf{number of
    instructions done per second}. A 3 GHz processor that needs 4 cycles per instruction runs
    $\frac{3 \times 10^9}{4} = 750$ million instructions per second. A 2 GHz processor that runs 2 instructions
    per cycle runs $4 \times 10^9$ instructions per second: it is faster, even with a lower frequency.
\end{intabih}

%================================================
\section{Memory}
%================================================

\subsection{What does memory look like inside?}
Imagine memory as a \textbf{big building full of boxes}. Each box has a number called the
\textbf{address}, and each box holds \textbf{one byte}.

\begin{center}
\begin{tikzpicture}[cellx/.style={draw=brandDark, fill=white, minimum width=2.4cm, minimum height=0.6cm, font=\codefont},
    addr/.style={font=\codefont\small, text=tealDeep}]
    \foreach \a/\v [count=\i from 0] in {0000/3A, 0001/FF, 0002/07, 0003/41, 0004/00} {
        \node[cellx] (c\i) at (0, -0.6*\i) {\v};
        \node[addr, left=0.2cm of c\i] {\a};
    }
    \node[font=\small\bfseries] at (-2.15, 0.6) {address};
    \node[font=\small\bfseries] at (0, 0.6) {content};
    \node[font=\small, text width=6.5cm, align=left] at (5.5,-1.2)
        {Address 0003 holds \hex{41}.\\
         The address never changes; the content can change.};
\end{tikzpicture}
\end{center}

\begin{tashbih}
    The address is like the \textbf{house number} in a street, and the content is like \textbf{the person
    who lives there}. The house number is fixed; the people living there can change.
\end{tashbih}

\subsection{Only two operations: read and write}
\begin{itemize}
    \item \textbf{Read}: the processor sends an address, and the memory gives back the content (the content does not change).
    \item \textbf{Write}: the processor sends an address and a value, and the memory stores the value instead of the old one.
\end{itemize}

\begin{mithal}
    Box 0002 holds \hex{07}. The processor runs ``write \hex{55} at 0002''. After that, the box holds
    \hex{55}, and \hex{07} is lost for ever. So we must be careful where we write: writing in the wrong place
    erases important data.
\end{mithal}

\begin{tarif}
    \term{Random Access Memory}{mémoire à accès aléatoire} (RAM): ``random access'' means we can reach
    \textbf{any box directly}, in the same time, without going through the boxes before it. This is the
    opposite of an old cassette tape, which you must wind until you reach the song.
\end{tarif}

\subsection{Types of memory}
\begin{center}
\begin{tikzpicture}[level distance=1.4cm,
    level 1/.style={sibling distance=7cm},
    level 2/.style={sibling distance=2.3cm},
    every node/.append style={unit, font=\small, minimum width=1.9cm},
    edge from parent/.style={draw, brandMuted, thick}]
    \node[unit=brandDark, fill=brandDark!10] {main memory}
        child {node[unit=brandOrange] {RAM\\{\scriptsize read and write}}
            child {node[unit=brandOrange] {SRAM\\{\scriptsize fast}}}
            child {node[unit=brandOrange] {DRAM\\{\scriptsize cheap}}}}
        child {node[unit=brandTeal] {ROM\\{\scriptsize read only}}
            child {node {PROM}}
            child {node {EPROM}}
            child {node {EEPROM\\{\scriptsize and Flash}}}};
\end{tikzpicture}
\end{center}

\begin{center}
\small
\begin{tabularx}{\linewidth}{|>{\centering\arraybackslash}p{2.2cm}|R|>{\centering\arraybackslash}p{2.4cm}|R|}
    \hline
    \ras{Type} & \ras{Description} & \ras{When power is off} & \ras{Used for} \\ \hline
    SRAM & transistors only, very fast and expensive & erased & cache memory inside the processor \\ \hline
    DRAM & a tiny capacitor per bit, must be ``refreshed'' all the time & erased & main memory (RAM sticks) \\ \hline
    ROM  & written once in the factory & kept & old boot programs \\ \hline
    EPROM & erased with ultraviolet light, then written again & kept & prototypes \\ \hline
    EEPROM / Flash & erased electrically & kept & BIOS/UEFI, USB keys, SSD \\ \hline
\end{tabularx}
\end{center}

\begin{tashbih}
    \textbf{SRAM} is like a light switch: it stays in its position (on or off) as long as there is power.
    \textbf{DRAM} is like a glass with a small hole: the water (the electric charge) leaks slowly, so we must
    fill it again thousands of times per second. This is called \textbf{refresh}. But the glass is much
    smaller than the switch, so we can put billions of them on a cheap memory stick.
\end{tashbih}

\begin{hal}
    Why does a computer need ROM or Flash? When you switch it on, the RAM is completely empty. The processor
    needs a first program to start with: the boot program (BIOS or UEFI), kept in a memory that is not erased.
    This program checks the hardware, then finds the operating system on the disk and loads it into RAM.
\end{hal}

\begin{tarif}
    A \term{volatile}{volatile} memory loses its content when the power is off (RAM).
    A \term{non-volatile}{non volatile} memory keeps it (ROM, Flash).
\end{tarif}

\subsection{Calculating the size of a memory}
If the processor has $n$ address lines (wires), it can write $2^n$ different addresses.

\begin{tarif}
    $\text{number of boxes} = 2^{n}$ \qquad and \qquad
    $\text{size} = 2^{n} \times \text{size of one box}$
\end{tarif}

\begin{khatwa}[: the memory size of the 8086]
    The 8086 has a \textbf{20-bit} address bus, and each box holds one byte.
    \begin{enumerate}
        \item Number of addresses: $2^{20}$.
        \item $2^{20} = 2^{10} \times 2^{10} = 1024 \times 1024 = 1\,048\,576$ addresses.
        \item Size $= 2^{20} \times 1$ byte $= 1$ megabyte (1 MB).
    \end{enumerate}
\end{khatwa}

\begin{khatwa}[: how many address lines for a 64 KB memory?]
    \begin{enumerate}
        \item $64\ \text{KB} = 64 \times 1024 = 2^{6} \times 2^{10} = 2^{16}$ bytes.
        \item So $2^{n} = 2^{16}$, and $n = 16$ lines.
    \end{enumerate}
\end{khatwa}

%================================================
\section{Buses}
%================================================

\begin{tarif}
    A \term{bus}{bus}: a group of parallel wires that carries information between the processor, the memory
    and the input/output devices. Each wire carries one bit at a time.
\end{tarif}

\begin{center}
\begin{tikzpicture}
    \node[unit, minimum width=2.8cm, minimum height=4cm] (cpu) at (-5,0) {Processor};
    \node[unit=brandPurple, minimum width=2.6cm, minimum height=1.2cm] (mem) at (0,2.6) {Memory};
    \node[unit=brandGreen, minimum width=2.6cm, minimum height=1.2cm] (io) at (4,2.6) {Input and\\output};
    \draw[brandOrange, line width=5pt] (-3.6,0.9) -- (6,0.9);
    \draw[brandTeal!80!black, line width=5pt] (-3.6,0) -- (6,0);
    \draw[brandMuted, line width=5pt] (-3.6,-0.9) -- (6,-0.9);
    \node[font=\small\bfseries, text=orangeDeep, anchor=west] at (6.1,0.9) {address bus};
    \node[font=\small\bfseries, text=tealDeep, anchor=west] at (6.1,0) {data bus};
    \node[font=\small\bfseries, text=brandMuted, anchor=west] at (6.1,-0.9) {control bus};
    \foreach \x in {0, 4} {
        \draw[brandOrange, line width=2pt, -{Stealth}] (\x-0.6,0.9) -- (\x-0.6,1.95);
        \draw[brandTeal!80!black, line width=2pt, {Stealth}-{Stealth}] (\x,0) -- (\x,1.95);
        \draw[brandMuted, line width=2pt, {Stealth}-{Stealth}] (\x+0.6,-0.9) -- (\x+0.6,1.95);
    }
\end{tikzpicture}
\end{center}

\begin{center}
\begin{tabularx}{\linewidth}{|>{\centering\arraybackslash}p{3cm}|>{\centering\arraybackslash}p{2.6cm}|R|}
    \hline
    \ras{Bus} & \ras{Direction} & \ras{What does it carry?} \\ \hline
    \textbf{address} bus & from the processor only & the number of the box the processor wants to read or write \\ \hline
    \textbf{data} bus & both ways & the value that is read or written \\ \hline
    \textbf{control} bus & both ways & signals such as read (RD), write (WR), clock, interrupt \\ \hline
\end{tabularx}
\end{center}

\begin{tashbih}
    Imagine you ask for a parcel at the post office. The \textbf{address} tells the clerk which shelf you mean
    (address bus). The \textbf{parcel} itself is the data (data bus). The words ``give me'' or ``take this''
    are the control signal (read or write).
\end{tashbih}

\begin{khatwa}[: the processor reads box 0003]
    \begin{enumerate}
        \item The processor puts the address 0003 on the \textbf{address bus}.
        \item It sends a \textbf{read} signal (RD) on the \textbf{control bus}.
        \item The memory puts the content of the box (\hex{41}) on the \textbf{data bus}.
        \item The processor takes the value \hex{41} and keeps it in one of its registers.
    \end{enumerate}
\end{khatwa}

\subsection{Bus width}
\begin{itemize}
    \item The \textbf{width of the address bus} decides the \textbf{maximum memory size}: $n$ lines give $2^n$ addresses.
    \item The \textbf{width of the data bus} decides \textbf{how many bits move at once}: a 16-bit bus moves two bytes at a time.
\end{itemize}

\begin{khatwa}[: moving 1000 bytes with different buses]
    \begin{enumerate}
        \item An \textbf{8-bit} data bus moves one byte at a time: we need 1000 transfers.
        \item A \textbf{16-bit} bus moves two bytes: $1000 \div 2 = 500$ transfers.
        \item A \textbf{64-bit} bus moves 8 bytes: $1000 \div 8 = 125$ transfers only.
    \end{enumerate}
    The wider the bus, the faster the data moves, just as more cars pass on an 8-lane motorway than on a one-lane road.
\end{khatwa}

\begin{center}
\begin{tabular}{|c|c|c|c|}
    \hline
    \ras{Processor} & \ras{Data bus} & \ras{Address bus} & \ras{Max memory} \\ \hline
    Intel 8086 & 16 bits & 20 bits & $2^{20}$ = 1 MB \\ \hline
    Intel 80386 & 32 bits & 32 bits & $2^{32}$ = 4 GB \\ \hline
    modern x86-64 & 64 bits & 48 bits (in practice) & $2^{48}$ = 256 TB \\ \hline
\end{tabular}
\end{center}

%================================================
\section{Exercises}
%================================================

\begin{tamrin}
    Convert these numbers: (a) \hex{3C} to decimal and binary. (b) The decimal number 200 to hexadecimal.
\end{tamrin}
\begin{hall}
    (a) $\hex{3C} = 3 \times 16 + 12 = 60$, and in binary $\bin{0011\,1100}$.\\
    (b) $200 = 12 \times 16 + 8$, and 12 is written C, so $200 = \hex{C8}$.
\end{hall}

\begin{tamrin}
    A processor has a 24-bit address bus, and each memory box holds one byte. What is the maximum memory size?
\end{tamrin}
\begin{hall}
    $2^{24} = 2^{4} \times 2^{20} = 16 \times 1\ \text{MB} = 16\ \text{MB}$.
\end{hall}

\begin{tamrin}
    We want a 4 GB memory. How many lines does the address bus need?
\end{tamrin}
\begin{hall}
    $4\ \text{GB} = 2^{2} \times 2^{30} = 2^{32}$ bytes, so we need \textbf{32 lines}.
\end{hall}

\begin{tamrin}
    True or false? Explain.
    (a) In the Von Neumann model, the program is kept in a different memory from the data.
    (b) The address bus carries information in both directions.
    (c) DRAM keeps its content after the computer is switched off.
\end{tamrin}
\begin{hall}
    (a) \textbf{False}: Von Neumann has one memory for program and data; separation is the Harvard model.\\
    (b) \textbf{False}: only the processor sends addresses, so it goes one way.\\
    (c) \textbf{False}: DRAM is volatile; it is erased when the power goes off.
\end{hall}

\begin{tamrin}
    A processor runs at 2 GHz, and each instruction needs 4 cycles on average. How many instructions does it run per second?
\end{tamrin}
\begin{hall}
    Cycles per second $= 2 \times 10^{9}$. Instructions $= \dfrac{2 \times 10^{9}}{4} = 0.5 \times 10^{9}$,
    that is \textbf{500 million instructions per second}.
\end{hall}
